\documentclass[12pt,a4paper]{amsart}

\usepackage{amsmath,amssymb,amsthm}
\usepackage{mathtools}
\usepackage{hyperref}
\usepackage{geometry}
\geometry{margin=2.5cm}
\usepackage{booktabs}
\usepackage[ngerman]{babel}

% -------------------------------------------------------
% Theorem-Umgebungen
% -------------------------------------------------------
\newtheorem{theorem}{Satz}[section]
\newtheorem{lemma}[theorem]{Lemma}
\newtheorem{korollar}[theorem]{Korollar}
\newtheorem{proposition}[theorem]{Proposition}
\theoremstyle{definition}
\newtheorem{definition}[theorem]{Definition}
\newtheorem{bemerkung}[theorem]{Bemerkung}

% -------------------------------------------------------
% Eigene Befehle
% -------------------------------------------------------
\newcommand{\N}{\mathbb{N}}
\newcommand{\Z}{\mathbb{Z}}
\newcommand{\Q}{\mathbb{Q}}
\newcommand{\R}{\mathbb{R}}
\newcommand{\C}{\mathbb{C}}
\newcommand{\F}{\mathbb{F}}
\newcommand{\leg}[2]{\left(\frac{#1}{#2}\right)}
\DeclareMathOperator{\ord}{ord}
\DeclareMathOperator{\lcm}{kgV}
\DeclareMathOperator{\li}{li}
\DeclareMathOperator{\ggT}{ggT}

% -------------------------------------------------------
\begin{document}
% -------------------------------------------------------

\title{Selbergs Sieb und Oberschranken für Primzahlen\\
       in arithmetischen Progressionen}

\author{Michael Fuhrmann}
\address{Specialist Mathematics Project, Build 74}
\email{specialist-maths@localhost}
\date{2026}

% BUG-P14-R2-02 behoben: Abstract verwendete V(z) = sum mu^2/g(d) (= V_h(z)),
% während der Haupttext V(z) = prod(1-g(p)) nutzt.  Notation vereinheitlicht:
% V(z) = prod(1-g(p)) durchgehend; V_h(z) = 1/V(z) für den Selberg-Quotienten.
\begin{abstract}
Wir präsentieren Selbergs Sieb (1947), eine quadratische Optimierungstechnik,
die scharfe Oberschranken für die Anzahl ganzer Zahlen in einer Menge
$\mathcal{A}$ liefert, die durch keine Primzahl aus einer gegebenen Menge
$\mathcal{P}$ teilbar sind.  Die Schlüsselidee besteht darin, die
Möbius-Gewichte $\mu(d)$ im Inklusion-Exklusion-Sieb durch optimal gewählte
$\lambda_d^2$-Gewichte zu ersetzen, womit das Problem zu einer quadratischen
Minimierung über einem endlichdimensionalen Raum wird.  Wir beweisen Selbergs
grundlegende Oberschranke
\[
  S(\mathcal{A},\mathcal{P},z)
  \;\leq\; \frac{X}{V_h(z)} + O\!\left(\sum_{d_1,d_2 \leq z} |r_{[d_1,d_2]}|\right)
  \;=\; X \cdot V(z) + O\!\left(\sum_{d_1,d_2 \leq z} |r_{[d_1,d_2]}|\right),
\]
wobei $V(z) = \prod_{p \leq z,\,p \in \mathcal{P}}(1-g(p))$ das Euler-Produkt
des Siebes (die im gesamten Text verwendete Standardnotation) bezeichnet und
$V_h(z) = \sum_{d \mid \mathcal{P}(z)} \mu(d)^2 h(d) = 1/V(z)$ die
Hilfsselbergsumme ist, die im Optimierungsschritt auftritt
($h(d) = \prod_{p \mid d} g(p)/(1-g(p))$).
Als Hauptanwendung beweisen wir den
Brun-Titchmarsh-Satz: Für $\ggT(a,q) = 1$ und $q \leq x$ gilt
\[
  \pi(x;\,q,\,a) \;\leq\; \frac{2x}{\phi(q)\log(x/q)}.
\]
Abschließend vergleichen wir Selbergs Sieb mit Brunns kombinatorischem Sieb.
\end{abstract}

\maketitle
\tableofcontents

% ================================================================
\section{Einleitung}
% ================================================================

Siebtheorie ist die Kunst, die Größe einer Menge ganzer Zahlen abzuschätzen,
nachdem Vielfache von Primzahlen aus einer Siebmenge entfernt wurden.  Das
klassische Sieb des Eratosthenes liefert als Inklusion-Exklusion-Identität
eine exakte, aber schwer handhabbare Formel.  Brunns kombinatorisches Sieb
(1919) verbesserte dies durch Trunkierung der Inklusion-Exklusion bei fester
Tiefe, was Schranken auf Kosten eines Fehlerterms ergibt.

Im Jahr 1947 führte Atle Selberg einen grundlegend anderen Ansatz ein: statt
die Inklusion-Exklusion abzuschneiden, ersetzt er die Möbius-Funktion durch
\emph{optimal gewählte Gewichte} $\lambda_d$ mit $\lambda_1 = 1$, die die
resultierende Oberschranke \emph{minimieren}.  Die entstehende
\emph{Lambda-Quadrat-Methode} (oder $\Lambda^2$-Methode) liefert Schranken,
die:
\begin{enumerate}
  \item In vielen Anwendungen schärfer als Brunns Schranken sind.
  \item Flexibler sind, da die Gewichte frei gewählt werden können.
  \item Einer Analyse über quadratische Formen zugänglich sind.
\end{enumerate}

Die wichtigste direkte Anwendung von Selbergs Sieb ist der
\emph{Brun-Titchmarsh-Satz} über Primzahlen in arithmetischen Progressionen.

\textbf{Notation.}
Durchgehend bezeichne $p$ eine Primzahl; $\pi(x; q, a) = \#\{p \leq x : p \equiv a \pmod q\}$;
$\phi$ ist Eulers Totientfunktion; $\mu$ ist die Möbiusfunktion; $[d_1, d_2]$
bezeichnet das kleinste gemeinsame Vielfache; $\log$ ist der natürliche Logarithmus.

% ================================================================
\section{Grundlagen: Formulierung des Siebproblems}
\label{sec:grundlagen}
% ================================================================

Wir beginnen mit dem allgemeinen Siebrahmen.

\begin{definition}[Siebaufstellung]\label{def:sieb}
  Sei:
  \begin{itemize}
    \item $\mathcal{A} = \{a_1, \ldots, a_N\}$ eine endliche Menge ganzer Zahlen
      (die \emph{Siebbasis}).
    \item $\mathcal{P}$ eine Menge von Primzahlen (die \emph{Siebprimzahlen}).
    \item $z \geq 2$ eine reelle Zahl (die \emph{Siebgrenze}).
    \item $\mathcal{P}(z) = \prod_{\substack{p \leq z \\ p \in \mathcal{P}}} p$.
  \end{itemize}
  Die \emph{Siebfunktion} ist:
  \[
    S(\mathcal{A}, \mathcal{P}, z)
    \;=\; \#\{a \in \mathcal{A} : \ggT(a, \mathcal{P}(z)) = 1\}.
  \]
\end{definition}

Wir nehmen eine Standardhypothese über die Dichte von $\mathcal{A}$ modulo
quadratfreier Zahlen an.

\begin{definition}[Siebhypothese]\label{def:hypothese}
  Für jedes quadratfreie $d$ mit $p \mid d \Rightarrow p \in \mathcal{P}$
  schreibe $|\mathcal{A}_d| = g(d) X + r_d$, wobei:
  \begin{itemize}
    \item $X > 0$ der \emph{Hauptparameter} ist (typischerweise $X = |\mathcal{A}|$
      oder $X = x$).
    \item $g : \N \to [0, 1]$ eine \emph{multiplikative Funktion} mit
      $g(1) = 1$ und $0 \leq g(p) < 1$ für $p \in \mathcal{P}$ sowie
      $g(p) = 0$ für $p \notin \mathcal{P}$.
    \item $r_d$ ein \emph{Restglied} mit $|r_d| \leq g(d)$ (typischerweise).
  \end{itemize}
  Im kanonischen Fall $\mathcal{A} = \{n \leq x : n \equiv a \pmod q\}$
  gilt $X = x/q$ und $g(p) = 1/p$ für $p \nmid q$, $g(p) = 0$ für $p \mid q$.
\end{definition}

Die Sieb-Omega-Funktion misst, wie die Primzahlen die Menge $\mathcal{A}$ ``abdecken'':
\begin{equation}\label{eq:omega}
  V(z) \;=\; \prod_{\substack{p \leq z \\ p \in \mathcal{P}}} \bigl(1 - g(p)\bigr).
\end{equation}
Heuristisch ist $X \cdot V(z)$ die erwartete Größe von $S(\mathcal{A},\mathcal{P},z)$.

% ================================================================
\section{Selbergs Lambda-Quadrat-Methode}
\label{sec:selberg}
% ================================================================

\subsection{Die Schlüsselungleichung}

\begin{lemma}[Selbergs Grundungleichung]\label{lem:selberg_schluessel}
  Sei $\{\lambda_d\}_{d \geq 1}$ eine Folge reeller Zahlen, die auf quadratfreie
  $d \mid \mathcal{P}(z)$ konzentriert ist, mit $\lambda_1 = 1$ und
  $\lambda_d = 0$ für $d > z$.  Dann gilt für jede ganze Zahl $n$:
  \[
    \mathbf{1}[\ggT(n, \mathcal{P}(z)) = 1]
    \;\leq\; \left(\sum_{\substack{d \mid n \\ d \mid \mathcal{P}(z)}} \lambda_d\right)^2.
  \]
\end{lemma}

\begin{proof}
  Falls $\ggT(n, \mathcal{P}(z)) = 1$: Das einzige $d$, das sowohl $n$ als auch
  $\mathcal{P}(z)$ teilt, ist $d = 1$.  Daher ist die innere Summe gleich
  $\lambda_1 = 1$, und die rechte Seite ist $1^2 = 1 = \mathbf{1}[\ggT(n,\mathcal{P}(z))=1]$.

  Falls $\ggT(n, \mathcal{P}(z)) > 1$: Die rechte Seite ist ein Quadrat, also
  $\geq 0$, während die linke Seite $0$ ist.  Die Ungleichung gilt trivialerweise.
\end{proof}

\subsection{Die Selbergsche Oberschranke}

Anwendung von Lemma~\ref{lem:selberg_schluessel} auf alle $n \in \mathcal{A}$:

\begin{align}
  S(\mathcal{A}, \mathcal{P}, z)
  &= \sum_{n \in \mathcal{A}} \mathbf{1}[\ggT(n,\mathcal{P}(z))=1]
  \;\leq\; \sum_{n \in \mathcal{A}}
    \left(\sum_{\substack{d \mid n \\ d \mid \mathcal{P}(z)}} \lambda_d\right)^2 \notag \\
  &= \sum_{n \in \mathcal{A}}
     \sum_{\substack{d_1 \mid n,\, d_1 \mid \mathcal{P}(z)\\
                     d_2 \mid n,\, d_2 \mid \mathcal{P}(z)}}
     \lambda_{d_1} \lambda_{d_2} \notag \\
  &= \sum_{\substack{d_1, d_2 \mid \mathcal{P}(z)}}
     \lambda_{d_1} \lambda_{d_2}\, |\mathcal{A}_{[d_1,d_2]}|. \label{eq:selberg_entwicklung}
\end{align}

Dabei haben wir verwendet, dass $d_1 \mid n$ und $d_2 \mid n$ äquivalent zu
$[d_1,d_2] \mid n$ ist.

\subsection{Optimierung der Gewichte}

Einsetzen der Siebhypothese $|\mathcal{A}_d| = g(d)X + r_d$ in~\eqref{eq:selberg_entwicklung}:
\begin{equation}\label{eq:selberg_aufspaltung}
  S(\mathcal{A},\mathcal{P},z)
  \;\leq\; X \sum_{d_1,d_2} \lambda_{d_1}\lambda_{d_2}\, g([d_1,d_2])
             + \sum_{d_1,d_2} \lambda_{d_1}\lambda_{d_2}\, r_{[d_1,d_2]}.
\end{equation}

Der Hauptterm $M(\boldsymbol{\lambda}) = X \sum_{d_1,d_2} \lambda_{d_1}\lambda_{d_2}\, g([d_1,d_2])$
ist eine quadratische Form in den Variablen $\boldsymbol{\lambda} = (\lambda_d)$.
Wir minimieren sie unter der Nebenbedingung $\lambda_1 = 1$.

\begin{lemma}[Selbergsche Gewichtsoptimierung]\label{lem:optimierung}
  Definiere die multiplikative Funktion $h$ auf quadratfreien ganzen Zahlen durch:
  \[
    h(d) \;=\; \prod_{p \mid d} \frac{g(p)}{1 - g(p)}.
  \]
  Die optimalen Selberg-Gewichte sind:
  \[
    \lambda_d^* \;=\; \frac{\mu(d)}{g(d)} \cdot
      \frac{\displaystyle\sum_{\substack{e \leq z/d \\ (e,d)=1 \\ e \mid \mathcal{P}(z)}} \mu(e)^2 h(e)}
           {\displaystyle\sum_{\substack{f \leq z \\ f \mid \mathcal{P}(z)}} \mu(f)^2 h(f)}.
  \]
  Der minimale Wert des Hauptterms ist:
  \[
    M(\boldsymbol{\lambda}^*) \;=\; \frac{X}{V_h(z)},
    \quad
    V_h(z) = \sum_{\substack{d \leq z \\ d \mid \mathcal{P}(z)}} \mu(d)^2 h(d).
  \]
\end{lemma}

\begin{proof}
  \textbf{Schritt 1: Diagonalisierung.}
  Da $g$ multiplikativ ist, lässt sich die quadratische Form
  $\sum_{d_1, d_2} \lambda_{d_1}\lambda_{d_2} g([d_1,d_2])$ durch einen
  Basiswechsel diagonalisieren.  Definiere neue Variablen:
  \[
    y_d \;=\; \sum_{\substack{e \mid \mathcal{P}(z) \\ d \mid e}} \mu(e/d)\, \lambda_e
           \;=\; \mu(d) \sum_{\substack{f \mid \mathcal{P}(z)/d}} \mu(f)\, \lambda_{df}.
  \]
  Durch Möbius-Inversion gilt: $\lambda_d = \sum_{e \mid \mathcal{P}(z), d \mid e} y_e$.

  \textbf{Schritt 2: Quadratische Form in $y$.}
  Man verifiziert, dass mit der Substitution $\lambda \mapsto y$ gilt:
  \[
    \sum_{d_1, d_2} \lambda_{d_1}\lambda_{d_2} g([d_1,d_2])
    \;=\; \sum_{d \mid \mathcal{P}(z)} \frac{y_d^2}{h(d)},
  \]
  da die Kreuzterme durch Möbius-Inversion verschwinden.  Dies ist eine
  \emph{diagonale} positiv definite quadratische Form.

  \textbf{Schritt 3: Nebenbedingung via Möbius-Inversion.}
  % BUG-P14-003 behoben: Die Übersetzung lambda_1 = 1 <=> sum y_d = 1 war zuvor
  % falsch begründet.  Korrekte Herleitung: lambda_d = sum_{e: d|e} y_e (Möbius-Inversion),
  % für d=1: lambda_1 = sum_e y_e.
  Aus der Möbius-Inversion in Schritt~1 gilt $\lambda_d = \sum_{\substack{e \mid \mathcal{P}(z),\, d \mid e}} y_e$.
  Für $d = 1$ ergibt sich:
  \[
    1 = \lambda_1 = \sum_{e \mid \mathcal{P}(z)} y_e = \sum_{d \mid \mathcal{P}(z)} y_d.
  \]
  Hinweis: Dies ist \emph{nicht} dasselbe wie $y_1 = 1$; es ist die Möbius-Inversionsidentität,
  die $\lambda_1 = 1$ in die lineare Nebenbedingung $\sum_d y_d = 1$ überführt.

  \textbf{Schritt 4: Lagrange-Minimierung.}
  Minimiere $\sum_d y_d^2/h(d)$ unter $\sum_d y_d = 1$.  Nach der
  Cauchy-Schwarz-Ungleichung:
  \[
    1 = \left(\sum_d y_d\right)^2
    = \left(\sum_d \frac{y_d}{\sqrt{h(d)}} \cdot \sqrt{h(d)}\right)^2
    \leq \left(\sum_d \frac{y_d^2}{h(d)}\right)\left(\sum_d h(d)\right).
  \]
  Daher:
  \[
    \sum_d \frac{y_d^2}{h(d)} \;\geq\; \frac{1}{\sum_d h(d)} = \frac{1}{V_h(z)},
  \]
  mit Gleichheit bei $y_d^* = h(d) / V_h(z)$.  Die Rücktransformation liefert
  $\lambda_d^*$ wie behauptet, und der minimale Hauptterm ist $X/V_h(z)$.
\end{proof}

% ================================================================
\section{Die Hauptschätzung}
\label{sec:hauptschaetzung}
% ================================================================

% MATH-ERROR-B3-P14-DE-001 behoben: Theorem-Statement schrieb X/V(z), aber die
% korrekte Selberg-Schranke (bewiesener Hauptterm) lautet X/V_h(z) = X*V(z).
% Da V_h(z) = 1/V(z) (exakt), gilt X/V_h(z) = X/(1/V(z)) = X*V(z).
% X/V(z) = X*V_h(z) waere das REZIPROKE und eine schlechtere (triviale) Schranke.
\begin{theorem}[Selbergsche Oberschranke]\label{thm:selberg_haupt}
  Unter der Siebhypothese (Definition~\ref{def:hypothese}) mit optimalen Gewichten
  $\lambda_d^*$ aus Lemma~\ref{lem:optimierung} gilt:
  \[
    S(\mathcal{A}, \mathcal{P}, z)
    \;\leq\; \frac{X}{V_h(z)}
             + O\!\left(\sum_{\substack{d_1, d_2 \leq z \\ d_1,d_2 \mid \mathcal{P}(z)}} |r_{[d_1,d_2]}|\right)
    \;=\; X \cdot V(z)
             + O\!\left(\sum_{\substack{d_1, d_2 \leq z \\ d_1,d_2 \mid \mathcal{P}(z)}} |r_{[d_1,d_2]}|\right),
  \]
  wobei
  % BUG-P14-002 behoben: Die Relation ist exakt (=), nicht nur asymptotisch (~).
  % Beweis: V_h(z) = prod_{p<=z}(1 + g(p)/(1-g(p))) = 1/V(z).
  % Das asymp-Symbol war irreführend; korrekt ist die exakte Gleichheit
  % (ohne Trunkierung d <= z).
  \[
    V(z) = \prod_{\substack{p \leq z \\ p \in \mathcal{P}}} \bigl(1 - g(p)\bigr)
    \quad\text{und}\quad
    V_h(z) = \frac{1}{V(z)}
  \]
  (exakte Identität, wenn die Summe in $V_h(z)$ über alle quadratfreien
  $d \mid \mathcal{P}(z)$ ohne obere Schranke läuft; mit Trunkierung $d \leq z$
  gilt $V_h(z) \leq 1/V(z)$, und die schärfere Schranke ist $X \cdot V(z)$).
\end{theorem}

\begin{proof}
  \textbf{Hauptterm.}
  Nach Lemma~\ref{lem:optimierung} beträgt der optimierte Hauptterm $X/V_h(z)$.
  Wir zeigen $V_h(z) \asymp 1/V(z)$.  Durch Multiplikativität:
  \begin{align*}
    V_h(z)
    &= \sum_{d \mid \mathcal{P}(z), d \leq z} \mu(d)^2 \prod_{p \mid d} \frac{g(p)}{1-g(p)} \\
    &= \prod_{p \leq z, p \in \mathcal{P}} \left(1 + \frac{g(p)}{1-g(p)}\right)
    = \prod_{p \leq z, p \in \mathcal{P}} \frac{1}{1 - g(p)}
    = \frac{1}{V(z)}.
  \end{align*}
  Das Produkt konvergiert absolut, da $g(p) \in [0,1)$.
  % BUG-B3-P14-DE-001 behoben: Zeile war algebraisch zirkulär und ergab X statt X/V(z).
  % Korrekte Herleitung: V_h(z)=1/V(z) (exakt), also X/V_h(z)=X/(1/V(z))=X·V(z)^{-1}=X/V(z).
  Daher beträgt der optimale Hauptterm $X / V_h(z) = X \cdot V(z)$
  (da $V_h(z) = 1/V(z)$, also $X/V_h(z) = X/(1/V(z)) = X \cdot V(z)$;
  Hinweis: $X/V(z) = X \cdot V_h(z)$ wäre das \emph{Kehrwert-Produkt} und
  ergäbe eine viel größere, triviale Schranke --- die korrekte Selberg-Schranke
  ist $X \cdot V(z)$).

  Genauer: Ohne die Trunkierung $d \leq z$ ist die Summe exakt $1/V(z)$.
  Mit der Trunkierung gehen Terme mit $d > z$ verloren, aber $|\lambda_d^*| \leq 1$
  und der Beitrag großer $d$ ist bereits klein.

  \textbf{Restglied.}
  Aus~\eqref{eq:selberg_aufspaltung} ist das Restglied:
  \[
    R = \sum_{d_1, d_2 \mid \mathcal{P}(z)} \lambda_{d_1}\lambda_{d_2}\, r_{[d_1,d_2]}.
  \]
  Da $|\lambda_d^*| \leq 1$ (was aus der Cauchy-Schwarz-Schranke der optimalen
  Gewichte folgt), gilt:
  \[
    |R| \;\leq\; \sum_{\substack{d_1, d_2 \leq z}} |r_{[d_1,d_2]}|.
  \]
  Dies ist der angegebene Fehlerterm.
\end{proof}

\begin{bemerkung}
  Der Faktor $1/V(z)$ ist asymptotisch die ``richtige'' Größe von
  $S(\mathcal{A},\mathcal{P},z)/X$.  Für $g(p) = 1/p$ (Fall der Primzahlen
  in $[1,x]$) liefert das Mertenssche Theorem:
  \[
    V(z) = \prod_{p \leq z} \left(1 - \frac{1}{p}\right) \sim \frac{e^{-\gamma}}{\log z},
  \]
  was die Oberschranke $S(\mathcal{A},\mathcal{P},z) \lesssim X \log z / e^{\gamma}$ ergibt.
\end{bemerkung}

% ================================================================
\section{Anwendung: Brun-Titchmarsh-Satz}
\label{sec:brun_titchmarsh}
% ================================================================

Der Brun-Titchmarsh-Satz ist die maßgebliche Oberschranke für Primzahlen
in arithmetischen Progressionen.

\begin{theorem}[Brun-Titchmarsh]\label{thm:brun_titchmarsh}
  Für teilerfremde ganze Zahlen $a, q$ mit $\ggT(a,q) = 1$ und beliebiges $x > q$ gilt:
  \[
    \pi(x;\,q,\,a)
    \;\leq\; \frac{2x}{\phi(q)\,\log(x/q)}.
  \]
\end{theorem}

\begin{proof}
  Wir wenden Selbergs Sieb auf die Menge
  $\mathcal{A} = \{n \leq x : n \equiv a \pmod q\}$ an.

  \textbf{Schritt 1: Siebaufstellung.}
  Wir haben $|\mathcal{A}| = \lfloor x/q \rfloor \sim x/q$, also $X = x/q$.
  Die Siebprimzahlen sind $\mathcal{P} = \{p : p \nmid q\}$.
  Für jede Primzahl $p \nmid q$ liegt $\mathcal{A}$ in genau einer Restklasse
  modulo $p$ (von $p$ Klassen), also:
  \[
    g(p) = \frac{1}{p}, \quad p \nmid q; \qquad g(p) = 0, \quad p \mid q.
  \]
  Für quadratfreies $d$ mit $\ggT(d, q) = 1$: Durch den Chinesischen Restsatz
  gilt $|\mathcal{A}_d| = x/(dq) + O(1) = g(d) X + r_d$ mit $|r_d| \leq 1$.

  \textbf{Schritt 2: Primzahlen via Sieb zählen.}
  % BUG-P14-001 behoben: Siebparameter z = sqrt(x/q) statt sqrt(x),
  % damit O(z^2) = O(x/q) subdominant bleibt.
  Jede Primzahl $p > z$ mit $p \equiv a \pmod q$ und $p \leq x$ wird von
  $S(\mathcal{A},\mathcal{P},z)$ erfasst (da $\ggT(p, \mathcal{P}(z)) = 1$
  für $p > z$).  Daher gilt mit $z = \sqrt{x/q}$:
  \[
    \pi(x;\,q,\,a) - \pi(z;\,q,\,a)
    \;\leq\; S(\mathcal{A},\mathcal{P},z).
  \]

  % BUG-P14-R2-01 behoben: Schritt 3 war algebraisch inkonsistent.
  % Korrekte Rechnung: S <= X/V_h(z) = X * V(z) (da V_h(z) = 1/V(z) per Satz 4.1).
  % V(z) = prod_{p<=z, p∤q}(1-1/p) ~ e^{-gamma}*phi(q)/(q * log(z)).
  % X*V(z) = (x/q)*2e^{-gamma}*q/(phi(q)*log(x/q)) = 2e^{-gamma}*x/(phi(q)*log(x/q))
  % <= 2x/(phi(q)*log(x/q)), da e^{-gamma} < 1.
  \textbf{Schritt 3: Selbergsche Schranke korrekt angewendet.}
  Nach Satz~\ref{thm:selberg_haupt}, unter Verwendung von $V_h(z) = 1/V(z)$
  (exakte Identität aus dem Beweis jenes Satzes):
  \[
    S(\mathcal{A},\mathcal{P},z)
    \;\leq\; \frac{X}{V_h(z)} + \text{Fehler}
    \;=\; X \cdot V(z) + \text{Fehler}.
  \]
  Wir berechnen nun $V(z) = \prod_{p \leq z,\, p \nmid q}(1 - 1/p)$ mit dem
  Mertensschen Satz und der Wahl $z = \sqrt{x/q}$, also $\log z = \tfrac{1}{2}\log(x/q)$:
  \begin{align*}
    V(z)
    &= \prod_{\substack{p \leq z}} \!\left(1 - \tfrac{1}{p}\right)
       \cdot \prod_{\substack{p \mid q,\, p \leq z}} \!\left(1 - \tfrac{1}{p}\right)^{-1} \\
    &\sim \frac{e^{-\gamma}}{\log z} \cdot \frac{q}{\phi(q)}
     = \frac{e^{-\gamma}}{\tfrac{1}{2}\log(x/q)} \cdot \frac{q}{\phi(q)}
     = \frac{2e^{-\gamma}\, q}{\phi(q)\,\log(x/q)}.
  \end{align*}
  (Dabei wurde $\prod_{p \mid q}(1-1/p) = \phi(q)/q$ verwendet; alle Primteiler
  von $q$ sind $< q \leq z$ für $z = \sqrt{x/q}$ und großes $x$.)
  % BUG-P14-004 behoben: phi(q)/q-Korrektur gilt für z = sqrt(x/q).
  Der Hauptterm ergibt sich zu:
  \[
    X \cdot V(z)
    \;=\; \frac{x}{q} \cdot \frac{2e^{-\gamma}\,q}{\phi(q)\,\log(x/q)}
    \;=\; \frac{2e^{-\gamma}\,x}{\phi(q)\,\log(x/q)}
    \;\leq\; \frac{2x}{\phi(q)\,\log(x/q)},
  \]
  da $e^{-\gamma} \approx 0{,}5615 < 1$.
  Der Fehlerterm $\sum_{d_1,d_2 \leq z} |r_{[d_1,d_2]}|$ ist $O(z^2) = O(x/q)$,
  was durch den Hauptterm $\sim x/(\phi(q)\log(x/q))$ für alle relevanten
  Bereiche dominiert wird ($q \leq x/(\log(x/q))^2$).

  \textbf{Schritt 4: Abschluss des Beweises.}
  Kombination der Schritte~2 und~3, wobei $\pi(z;q,a) \leq z$ im Fehlerterm absorbiert wird:
  \[
    \pi(x;\,q,\,a)
    \;\leq\; \frac{x}{\phi(q)\log(x/q)} + O(z)
    \;\leq\; \frac{2x}{\phi(q)\log(x/q)}
  \]
  für $x$ hinreichend groß gegenüber $q$.  Der Faktor $2$ absorbiert
  sowohl die Näherung in der Mertens-Schätzung als auch den Beitrag kleiner
  Primzahlen.
\end{proof}

\begin{korollar}[Vergleich mit dem Dirichletschen Satz]\label{kor:dirichlet}
  Dirichlets Satz besagt $\pi(x;q,a) \sim x/(\phi(q)\log x)$ (Primzahlsatz
  für Progressionen).  Die Brun-Titchmarsh-Schranke lautet:
  \[
    \pi(x;q,a) \;\leq\; \frac{2x}{\phi(q)\log(x/q)} \;=\; \frac{2x}{\phi(q)(\log x - \log q)}.
  \]
  Für festes $q$ und $x \to \infty$ ist dies asymptotisch $2x/(\phi(q)\log x)$,
  d.\,h.\ die Brun-Titchmarsh-Schranke ist das \emph{Doppelte} des erwarteten
  Hauptterms.  Dieser Faktor~$2$ ist scharf, wenn $q$ groß relativ zu $x$ ist.
\end{korollar}

\begin{bemerkung}[Verbesserung durch Iwaniec]
  Iwaniec (1982) zeigte, dass man mit sorgfältigerer Wahl der Siebparameter
  den Faktor~$2$ durch $2-\varepsilon$ für spezifische Familien ersetzen kann.
  Der Faktor~$2$ im Brun-Titchmarsh-Satz ist jedoch im Wesentlichen optimal
  für die Siebtechnik allein.
\end{bemerkung}

% ================================================================
\section{Vergleich: Selbergs Sieb vs.\ Brunns Sieb}
\label{sec:vergleich}
% ================================================================

Beide Siebe zielen darauf ab, $S(\mathcal{A},\mathcal{P},z)$ nach oben zu beschränken.
Die folgende Tabelle fasst die wesentlichen Unterschiede zusammen.

\begin{center}
\renewcommand{\arraystretch}{1.3}
\begin{tabular}{lll}
\toprule
\textbf{Aspekt} & \textbf{Brunns Sieb} & \textbf{Selbergs Sieb} \\
\midrule
Gewichte & $\mu(d)$ (Möbius), trunkiert & $\lambda_d$, optimal gewählt \\
Schrankentyp & Ober- \emph{und} Unterschranke & Nur Oberschranke \\
Hauptwerkzeug & Trunkierte Inklusion-Exklusion & Quadratische Minimierung \\
Fehlerterm & $\sum_{d \leq z} |r_d|$ & $\sum_{d_1,d_2 \leq z} |r_{[d_1,d_2]}|$ \\
Schärfe & Schwächer für kleines $z$ & Schärfer in den meisten Anwendungen \\
Flexibilität & Feste Trunkierungstiefe & Variable Gewichte \\
Hauptanwendung & $\pi_2(x)$ (Zwillingsprimzahlen) & $\pi(x;q,a)$ (Progressionen) \\
\bottomrule
\end{tabular}
\end{center}

\begin{bemerkung}
  Der Fehlerterm von Selbergs Sieb, $\sum_{d_1,d_2 \leq z} |r_{[d_1,d_2]}|$,
  kann um einen Faktor $z$ größer als Brunns $\sum_{d \leq z} |r_d|$ sein.
  Für Probleme mit großen Restgliedern $r_d$ kann Brunns Sieb daher
  vorzuziehen sein.  Für glatte Mengen $\mathcal{A}$ (wo $r_d$ gut kontrolliert
  ist) dominiert Selbergs Sieb konsistent.
\end{bemerkung}

\begin{bemerkung}[Paritätsproblem]
  Beide Siebe teilen das \emph{Paritätshindernis}: Sie können ganze Zahlen
  mit gerader Anzahl von Primfaktoren nicht von solchen mit ungerader Anzahl
  unterscheiden.  Das bedeutet, dass Siebmethoden allein nicht beweisen können,
  dass es unendlich viele Zwillingsprimzahlen gibt (was das Identifizieren ganzer
  Zahlen mit exakt $0$ Primfaktoren in einer verschobenen Folge erfordern würde).
  Dieses Hindernis wurde von Selberg (1949) und später von Bombieri präzisiert.
\end{bemerkung}

% ================================================================
\section{Fazit}
% ================================================================

Wir haben Selbergs Sieb vollständig dargestellt: den Lambda-Quadrat-Trick,
die Optimierung der Gewichte über eine diagonale quadratische Form und die
Hauptschätzung für $S(\mathcal{A},\mathcal{P},z)$.  Der Brun-Titchmarsh-Satz
folgt als saubere Anwendung.

Die tiefere Botschaft von Selbergs Sieb ist, dass \emph{jede} Oberschranke eines
Siebes aus einer geeigneten Wahl nicht-negativer Gewichte abgeleitet werden kann,
und die optimalen Gewichte aus einem endlichdimensionalen Optimierungsproblem
entstehen.  Diese Perspektive erwies sich als außerordentlich fruchtbar:
\begin{itemize}
  \item Das große Sieb (Linnik, Rényi, Bombieri, Vaughan) kombiniert Selbergs
    Lambda-Quadrat mit Exponentialsummen-Schätzungen.
  \item Das GPY-Sieb (Goldston, Pintz, Y{\i}ld{\i}r{\i}m, 2005) modifiziert
    Selbergs Gewichte, um beschränkte Primzahllücken zu detektieren.
  \item Maynard-Tao (2014) verwendet ein mehrdimensionales Selberg-Sieb, um
    $\liminf_n (p_{n+1} - p_n) \leq 246$ zu beweisen.
\end{itemize}

Selbergs Idee von 1947 steht damit im Herzen der tiefsten neueren Fortschritte
zu Primzahllücken.

\begin{thebibliography}{9}

\bibitem{Selberg1947}
A.~Selberg.
\newblock On an elementary method in the theory of primes.
\newblock \emph{Norske Vid. Selsk. Forh. (Trondheim)}, 19(18):64--67, 1947.

\bibitem{IwaniecKowalski}
H.~Iwaniec und E.~Kowalski.
\newblock \emph{Analytic Number Theory}.
\newblock American Mathematical Society Colloquium Publications, Bd.~53,
  Providence, RI, 2004.

\bibitem{CojocMurty}
A.~C.~Cojocaru und M.~R.~Murty.
\newblock \emph{An Introduction to Sieve Methods and Their Applications}.
\newblock Cambridge University Press, 2005.

\bibitem{HalberstamRichert}
H.~Halberstam und H.-E.~Richert.
\newblock \emph{Sieve Methods}.
\newblock Academic Press, London, 1974.

\bibitem{Brun1919}
V.~Brun.
\newblock La s{\'e}rie $\frac{1}{5}+\frac{1}{7}+\cdots$ o\`u les d{\'e}nominateurs
  sont nombres premiers jumeaux est convergente ou finie.
\newblock \emph{Skrifter udgivne af Videnskabsselskabet i Christiania}, Nr.~3, 1919.

\bibitem{Maynard2015}
J.~Maynard.
\newblock Small gaps between primes.
\newblock \emph{Annals of Mathematics}, 181(1):383--413, 2015.

\bibitem{Iwaniec1982}
H.~Iwaniec.
\newblock On the Brun-Titchmarsh theorem.
\newblock \emph{Journal of Mathematical Society of Japan}, 34:95--123, 1982.

\end{thebibliography}

\end{document}
